\chapter{Numbers, Algebra and Functions}
\label{ch:numbers}

\begin{tcolorbox}[feynbox={Sky}{BBg},title={Before and After}]
\textbf{Before this chapter you need:} arithmetic with whole numbers ---
adding, subtracting, multiplying and dividing. Nothing else.\\
\textbf{After it you will understand:} what a variable is, how to rearrange an
equation without breaking it, and what people mean when they call something a
\emph{function}. Every later chapter is written in the language of functions,
so this is the vocabulary the rest of the book assumes.
\end{tcolorbox}

This chapter contains no finance and almost no notation. Its job is to make sure
that when Chapter~\ref{ch:calculus1} writes $f(x)$ and asks how fast it changes,
the symbols are already familiar rather than a second obstacle stacked on top of
the idea.

%% ─────────────────────────────────────────────────────────────────────────────
\section{Fractions, Ratios and Percentages}
\label{sec:percentages}

\situation{
A share costs \$40 on Monday. On Tuesday it costs \$46. You want to describe the
change to someone who does not know the starting price. Saying ``it went up \$6''
is useless to them --- \$6 on a \$40 share is a large move, \$6 on a \$4{,}000
share is a rounding error. What they want is the change \emph{relative to where
it started}.
}

\problem{
Absolute differences cannot be compared across things of different size. A
\$6 move in Brazil and a \$6 move in China tell you nothing about which market
had the bigger day. To compare, every quantity must first be put on a common
scale, and the natural common scale is ``fraction of the starting value''.
}

\workedarith{
Start 40, end 46.

\medskip
\textbf{Change:} $46 - 40 = 6$.

\textbf{As a fraction of the start:} $\dfrac{6}{40} = \dfrac{3}{20} = 0.15$.

\textbf{As a percentage:} $0.15 \times 100 = 15\%$.

\medskip
Check it the other way round. If 40 grows by 15\%, the new value is
$40 \times (1 + 0.15) = 40 \times 1.15 = 46$. \checkmark

\medskip
Now the same move on a bigger share. Start 4{,}000, end 4{,}006:
$\dfrac{6}{4000} = 0.0015 = 0.15\%$. Same \$6, one hundredth of the move.
}

\formaldef{
For a quantity that starts at $P_0$ and ends at $P_1$, the \textbf{simple return}
is
\[
  R = \frac{P_1 - P_0}{P_0} = \frac{P_1}{P_0} - 1 .
\]
A \textbf{percentage} is a fraction multiplied by 100; the symbol \% means
``divide by 100''. The quantity $P_1/P_0$ is called the \textbf{gross return} or
\textbf{price relative}: a gross return of $1.15$ and a simple return of $0.15$
say the same thing.
}

\analogybreak{
``15\% up'' and ``15\% down'' are not opposite moves. Going up 15\% from 40 gives
46; coming back down 15\% from 46 gives $46 \times 0.85 = 39.1$, not 40. The
percentage is always taken relative to \emph{wherever you currently are}, and
after the first move that is a different place. This asymmetry is not a curiosity
--- it is the reason Chapter~\ref{ch:explog} replaces percentages with logarithms
for everything in this thesis.
}

\whyhere{
Every number in this thesis is a return, not a price. The five BRICS markets
trade at wildly different index levels, and only by converting to returns can
BOVESPA and SHANGHAI be placed on the same axis, fed to the same model, and
scored by the same metric.
}

%% ─────────────────────────────────────────────────────────────────────────────
\section{Variables and Rearranging Equations}
\label{sec:variables}

\situation{
A recipe says: to find the cost of a taxi, take \$3 for the pickup and add \$2
for every kilometre. You could write this out in words every time, or you could
write $C = 3 + 2k$ and let $k$ stand for however many kilometres you happen to
travel. The letter is shorthand for ``a number I have not decided yet''.
}

\problem{
Once the relationship is written down, you often need to run it backwards. The
formula gives cost from distance, but the question you actually have is ``I have
\$15 --- how far can I go?''. That means solving for $k$ instead of $C$, and doing
that safely requires a rule for what you are allowed to do to an equation.
}

\workedarith{
The rule is: \emph{whatever you do to one side, do to the other}. An equation is
a claim that two things are equal, and equals stay equal if treated identically.

\medskip
Start with $C = 3 + 2k$ and set $C = 15$:
\[
  15 = 3 + 2k
\]
Subtract 3 from both sides:
\[
  15 - 3 = 3 + 2k - 3 \quad\Longrightarrow\quad 12 = 2k
\]
Divide both sides by 2:
\[
  \frac{12}{2} = \frac{2k}{2} \quad\Longrightarrow\quad 6 = k
\]

\medskip
\textbf{Check by substitution:} $C = 3 + 2 \times 6 = 3 + 12 = 15$. \checkmark

\medskip
Substituting the answer back into the original equation is the only check you
ever need, and it costs ten seconds.
}

\formaldef{
A \textbf{variable} is a symbol standing for a number that may take different
values. A \textbf{constant} is a fixed number. An \textbf{equation} asserts that
two expressions are equal; to \textbf{solve} it for a variable is to produce an
equivalent equation with that variable alone on one side.

The operations that preserve equality are: adding or subtracting the same
quantity from both sides, and multiplying or dividing both sides by the same
\emph{non-zero} quantity.
}

\analogybreak{
The ``do the same to both sides'' rule has one trap: multiplying both sides by
zero. From $5 = 5$ we may multiply both sides by 0 and get $0 = 0$, which is
true but has lost all information --- and the step cannot be undone. This is why
division by zero is forbidden: it is the reverse of a step that destroys
information. The same caution reappears in Chapter~\ref{ch:linalg}, where a
matrix with determinant zero cannot be inverted for exactly this reason.
}

\whyhere{
Rearranging is how the definition of a log return, $r = \ln(P_1/P_0)$, is turned
into the reconstruction $P_1 = P_0 e^{r}$ used whenever a synthetic return series
must be turned back into a synthetic price path.
}

%% ─────────────────────────────────────────────────────────────────────────────
\section{What a Function Is}
\label{sec:functions}

\situation{
A vending machine has a keypad. You press A4 and a specific bar of chocolate
comes out. You press A4 again tomorrow and the \emph{same} kind of bar comes out.
Press B2 and something else comes out. The machine is a rule that turns each
input into exactly one output, reliably, every time.
}

\problem{
The reliability is the whole point. A machine that sometimes gave chocolate and
sometimes gave crisps for the same button would be useless --- you could not
reason about it or plan around it. Mathematics needs the same guarantee before
it can say anything about how outputs change when inputs change.
}

\workedarith{
Take the rule ``double it, then add one'', written $f(x) = 2x + 1$.

\medskip
\begin{tabular}{@{}llll@{}}
\toprule
Input $x$ & Working & Output $f(x)$ \\
\midrule
$0$ & $2 \times 0 + 1$ & $1$ \\
$1$ & $2 \times 1 + 1$ & $3$ \\
$2$ & $2 \times 2 + 1$ & $5$ \\
$3$ & $2 \times 3 + 1$ & $7$ \\
$-1$ & $2 \times (-1) + 1$ & $-1$ \\
\bottomrule
\end{tabular}

\medskip
Each input appears once and produces one output. Note that outputs \emph{may}
repeat for different inputs without breaking anything: for $g(x) = x^2$ we get
$g(2) = 4$ and $g(-2) = 4$. Two doors, same room. That is allowed. What is not
allowed is one door opening onto two rooms.
}

\formaldef{
A \textbf{function} $f$ from a set $A$ to a set $B$ assigns to each element
$x \in A$ exactly one element $f(x) \in B$. We write $f : A \to B$, read
``$f$ maps $A$ to $B$''. The symbol $\in$ means ``is an element of''.

The value $f(x)$ is the \textbf{image} of $x$. The notation is deliberately
mechanical: $f$ is the machine, $x$ is what goes in, $f(x)$ is what comes out.
}

\analogybreak{
The vending-machine picture suggests the machine is simple and you could open it
up and see the mechanism. Many functions in this thesis are not like that. A
trained neural network is a function --- one input gives one output, reliably ---
but it has hundreds of thousands of internal numbers and no human-readable
mechanism. ``Function'' guarantees only the input--output discipline, never
simplicity or interpretability.
}

\whyhere{
A generator in a GAN is a function: it takes a vector of random noise and returns
a synthetic return series. A discriminator is a function: it takes a series and
returns a score. Training changes \emph{which} function these are, but they are
functions at every moment, which is what makes it meaningful to ask how their
outputs respond to their inputs.
}

%% ─────────────────────────────────────────────────────────────────────────────
\section{Domain and Range}
\label{sec:domain}

\situation{
The vending machine has buttons A1 to D5 and nothing else. Pressing a letter that
is not on the keypad does nothing --- the question ``what does the machine give
for Z9?'' has no answer, because Z9 was never a legal input.
}

\problem{
Ignoring which inputs are legal produces confident nonsense. A formula that
divides by $x$ is silent about what happens at $x = 0$ until someone asks; a
formula involving $\sqrt{x}$ has nothing to say about $x = -1$ if we are working
with ordinary numbers. Stating the legal inputs up front prevents the error.
}

\workedarith{
\textbf{Example 1:} $f(x) = 1/x$.

Legal inputs: every number except $0$. At $x = 2$, $f = 0.5$; at $x = 0.1$,
$f = 10$; at $x = 0.001$, $f = 1000$. As $x$ shrinks toward zero the output grows
without bound, and at $x = 0$ there is no output at all.

\medskip
\textbf{Example 2:} $f(x) = \sqrt{x}$ over ordinary (real) numbers.

Legal inputs: $x \geq 0$. Outputs: $\sqrt{0} = 0$, $\sqrt{9} = 3$,
$\sqrt{2} \approx 1.414$. Every output is $\geq 0$, so although the legal inputs
run from $0$ upwards, so do the outputs.

\medskip
\textbf{Example 3:} $f(x) = x^2$.

Legal inputs: all numbers. Outputs: never negative, because a number times itself
is never negative. So the outputs occupy only half the number line even though
the inputs occupy all of it.
}

\formaldef{
The \textbf{domain} of $f$ is the set of inputs on which $f$ is defined. The
\textbf{codomain} is the set the outputs are declared to live in. The
\textbf{range} (or image) is the set of values $f$ actually attains.

For $f(x) = x^2$ with domain $\R$ (all real numbers), the codomain may be
declared as $\R$, but the range is only $[0, \infty)$ --- the non-negative
numbers. Range is a subset of codomain, sometimes a strict one.
}

\analogybreak{
Domain is often presented as a technicality to be recited and forgotten. It is
not. In Chapter~\ref{ch:explog} the domain of $\ln$ is exactly the strictly
positive numbers, and that single restriction is what forces prices --- never
returns --- to be the thing we take logarithms of. The restriction carries real
content about what the mathematics is willing to model.
}

\whyhere{
The generators in this thesis end with a $\tanh$ activation, whose range is
$(-1, 1)$. That is a deliberate choice of range: it guarantees the network cannot
emit an arbitrarily large value, and it is why the training data must first be
squashed into the same interval before the two can be compared.
}

%% ─────────────────────────────────────────────────────────────────────────────
\section{Linear Functions and Slope}
\label{sec:linear}

\situation{
You are walking up a straight ramp. For every 4 metres you move forward, you rise
1 metre. The ramp does not get steeper or shallower --- the trade-off between
forward and upward is the same everywhere on it. One number, $1/4$, describes the
entire ramp.
}

\problem{
Most relationships are not ramps. But a great deal of mathematics consists of
finding the ramp that best matches a curve near a point of interest, because
ramps are the only shape we can reason about effortlessly. Before that programme
can start, ``steepness'' has to be pinned to a number.
}

\workedarith{
Take $f(x) = 3x + 2$.

\medskip
\begin{tabular}{@{}llll@{}}
\toprule
$x$ & $f(x)$ & Change in $x$ & Change in $f$ \\
\midrule
$1$ & $5$  & --- & --- \\
$2$ & $8$  & $1$ & $3$ \\
$4$ & $14$ & $2$ & $6$ \\
$7$ & $23$ & $3$ & $9$ \\
\bottomrule
\end{tabular}

\medskip
From $x = 1$ to $x = 4$: heights 5 and 14, change $9$, horizontal distance $3$,
ratio $9/3 = 3$.

From $x = 4$ to $x = 7$: heights 14 and 23, change $9$, distance $3$,
ratio $3$.

\medskip
The ratio is 3 no matter which two points we pick. That constancy is precisely
what makes the function linear, and the 3 is exactly the number multiplying $x$
in the formula. The $+2$ shifts the whole line up without tilting it: at $x = 0$
the value is 2.
}

\formaldef{
A \textbf{linear function} (strictly, an \emph{affine} one) has the form
\[
  f(x) = mx + c ,
\]
where $m$ is the \textbf{slope} or \textbf{gradient} and $c$ is the
\textbf{intercept}, the value at $x = 0$. For any two distinct points
$(x_1, y_1)$ and $(x_2, y_2)$ on the line,
\[
  m = \frac{y_2 - y_1}{x_2 - x_1} ,
\]
often read ``rise over run''. Positive $m$ tilts upward, negative $m$ downward,
and $m = 0$ is flat.
}

\analogybreak{
Slope as ``rise over run'' silently assumes the line is not vertical. A vertical
line has zero run, and the ratio would require dividing by zero --- so vertical
lines simply have no slope, rather than an infinite one. Refusing to write
$m = \infty$ here is the same discipline that will matter in
Chapter~\ref{ch:heavytails}, where a distribution with an undefined fourth
moment is genuinely undefined rather than merely large.
}

\whyhere{
The derivative in Chapter~\ref{ch:calculus1} is nothing more than this slope,
computed over a vanishingly small run. Gradient descent --- the algorithm that
trains every model in this thesis --- reads that slope and steps downhill.
}

%% ─────────────────────────────────────────────────────────────────────────────
\section{Quadratics and Curvature}
\label{sec:quadratics}

\situation{
Throw a ball straight up. It rises, slows, stops, and falls. Its height against
time is not a ramp: the steepness changes continuously, from strongly positive at
the throw to zero at the top to strongly negative on the way down.
}

\problem{
A single slope cannot describe this path, because the path has a different slope
at every instant. We need a shape with built-in bend, and the simplest such shape
is the one produced by squaring.
}

\workedarith{
Take $f(x) = x^2$.

\medskip
\begin{tabular}{@{}lllll@{}}
\toprule
$x$ & $f(x)$ & Interval & Change in $f$ & Average slope \\
\midrule
$0$ & $0$  & --- & --- & --- \\
$1$ & $1$  & $0 \to 1$ & $1$ & $1$ \\
$2$ & $4$  & $1 \to 2$ & $3$ & $3$ \\
$3$ & $9$  & $2 \to 3$ & $5$ & $5$ \\
$4$ & $16$ & $3 \to 4$ & $7$ & $7$ \\
\bottomrule
\end{tabular}

\medskip
The average slope over each unit step is $1, 3, 5, 7$ --- rising by 2 every time.
The function is not merely increasing; the \emph{rate} at which it increases is
itself increasing at a steady rate. That steady second-level change is what
``curvature'' means, and here it equals 2.

\medskip
Note also the symmetry: $f(-3) = 9 = f(3)$. The curve is a mirror image about the
vertical line $x = 0$, which sits at the bottom of the bowl.
}

\formaldef{
A \textbf{quadratic function} has the form
\[
  f(x) = ax^2 + bx + c , \qquad a \neq 0 .
\]
Its graph is a \textbf{parabola}. If $a > 0$ it opens upward and has a single
lowest point; if $a < 0$ it opens downward and has a single highest point. That
turning point sits at $x = -b/(2a)$.

For $f(x) = x^2$ we have $a = 1$, $b = 0$, $c = 0$, so the turning point is at
$x = 0$, matching the table above.
}

\analogybreak{
A parabola has exactly one turning point, which makes ``find the lowest value''
easy: there is nowhere else it could be. Almost none of the functions minimised
in this thesis are like that. A neural network's loss surface has enormous
numbers of local dips, and an algorithm that walks downhill will stop at whichever
one it happens to reach. The single-minimum picture is the exception, not the rule
--- Chapter~\ref{ch:calculus2} returns to this.
}

\whyhere{
Squared error --- the quantity minimised when the TimeGAN autoencoder learns to
reconstruct its input --- is a quadratic in the reconstruction error. Its bowl
shape is why the reconstruction loss has a well-defined best answer, and why a
reported $R^2$ of $+0.995$ means the bottom of that bowl was very nearly reached.
}

%% ─────────────────────────────────────────────────────────────────────────────
\section{Composition of Functions}
\label{sec:composition}

\situation{
A factory line has two stations. The first paints the part; the second boxes it.
A raw part enters, a painted part passes between them, a boxed painted part comes
out. Nobody inspects the middle stage --- from outside, the line is one machine
turning raw parts into finished goods.
}

\problem{
Deep networks are precisely this: dozens of simple stations in a row. To train
them we must answer ``if I adjust station one slightly, how much does the final
output move?'' --- a question about the whole line, asked at a single station.
Answering it requires first being clear about what stacking functions even means.
}

\workedarith{
Let $g(x) = x + 3$ (the painter) and $f(u) = 2u$ (the boxer). Apply $g$ first,
then $f$.

\medskip
\begin{tabular}{@{}llll@{}}
\toprule
Input $x$ & After $g$: $x+3$ & After $f$: $2(x+3)$ \\
\midrule
$1$ & $4$ & $8$ \\
$2$ & $5$ & $10$ \\
$5$ & $8$ & $16$ \\
\bottomrule
\end{tabular}

\medskip
The combined rule is $f(g(x)) = 2(x + 3) = 2x + 6$. Check at $x = 5$:
$2 \times 5 + 6 = 16$. \checkmark

\medskip
\textbf{Order matters.} Reverse the stations --- box first, then paint:
$g(f(x)) = 2x + 3$. At $x = 5$ that gives $13$, not $16$. Painting a boxed part
is not the same as boxing a painted one.
}

\formaldef{
Given $g : A \to B$ and $f : B \to C$, the \textbf{composition} $f \circ g$ is the
function $A \to C$ defined by
\[
  (f \circ g)(x) = f\bigl(g(x)\bigr) .
\]
It is read ``$f$ after $g$'': the function written on the \emph{left} is applied
\emph{last}. Composition requires the range of $g$ to lie inside the domain of
$f$ --- the parts leaving the first station must be things the second station can
accept.
}

\analogybreak{
The factory picture suggests each station is independent, so you could swap or
remove one freely. In a trained network that is false. Each layer's weights have
been fitted assuming exactly what the previous layer sends it; extract a middle
layer and it produces nothing meaningful on raw input. The stations are separable
in notation only.
}

\whyhere{
A neural network is one long composition, and the chain rule of
Chapter~\ref{ch:calculus2} is the tool that differentiates it. Everything called
``backpropagation'' is the chain rule applied to a composition of this kind. The
TimeGAN recovery function composed with its embedder is a composition whose
round-trip fidelity we measure directly --- and require to be accurate to within
$10^{-3}$ before training is allowed to proceed.
}
